%! TEX root = ../bcd-19-20.tex

\chapter{Funcții implicite}


Funcțiile implicite sînt definite de ecuații care, în general, nu se pot
rezolva sau se rezolvă foarte dificil. De exemplu, dacă avem o ecuație
de forma:
\[
  xy + 2x\sin y = 0
\]
și vrem să exprimăm de aici funcția $ y = y(x) $, constatăm că acest
lucru nu este posibil, deoarece ecuația nu poate fi rezolvată pentru $ y $.
Astfel, vom spune în acest caz că $ y $ a fost definită \emph{implicit}
de ecuația de mai sus.
\index{funcții!implicite}

Trecem direct la exemplificarea noțiunilor teoretice pe exerciții rezolvate.

\vspace{0.5cm}

1. Funcția $ z = z(x, y) $ este definită implicit de ecuația:
\[
  (y + z) \sin z - y (x + z) = 0.
\]
Calculați expresia:
\[
  E = z \sin z \frac{\partial z}{\partial x} - y^2 \frac{\partial z}{\partial y}.
\]

\emph{Soluție:} Notăm expresia implicită dată cu $ F(x, y, z) $. Pentru a putea
exprima $ z = z(x, y) $, deci ca funcție, este necesar ca expresia $ F $ să depindă
\emph{funcțional} de $ z $, adică să nu aibă pe $ z $ doar ca parametru ori drept
constantă. Așadar, avem o \emph{condiție de existență} a funcției implicite
$ z = z(x, y) $, anume $ \dfrac{\partial F}{\partial z} \neq 0 $.

Presupunem acum că ne aflăm în ipoteza de mai sus, i.e.\ condiția de existență
este satisfăcută. Atunci, pentru a exprima $ \dfrac{\partial z}{\partial x} $
și $ \dfrac{\partial z}{\partial y} $, vom deriva expresia $ F $, deoarece
doar acolo apare funcția $ z $. Obținem:
\[
  \frac{\partial F}{\partial x} = \frac{\partial z}{\partial x} \sin z + %
  (y + z) \cos z \dfrac{\partial z}{\partial x} - %
  y\Big(1 + \dfrac{\partial z}{\partial x} \Big) = 0
\]
Rezultă:
\[
  \frac{\partial z}{\partial x} = \frac{z}{\sin z + (y + z) \cos z - y}.
\]

Similar:
\[
  \frac{\partial F}{\partial y} = \Big(1 + \frac{\partial z}{\partial y} \Big) \sin z + %
  (y + z) \cos z \frac{\partial z}{\partial y} - (x + z) - y \frac{\partial z}{\partial y} = 0.
\]
Rezultă:
\[
  \frac{\partial z}{\partial y} = \dfrac{x + z - \sin z}{\sin z + (y + z)\cos z - y}.
\]
Înlocuind în expresia cerută, obținem $ E = 0 $.

\vspace{0.5cm}

Un alt tip de exercițiu de care sîntem interesați este acela al extremelor pentru
funcții definite implicit.

2. Să se determine extremele funcției $ y = y(x) $, definită implicit de ecuația:
\[
  x^3 + y^3 - 2xy = 0.
\]

\emph{Soluție:} Fie $ F(x, y) = x^3 + y^3 - 2xy $, astfel încît condiția dată
este $ F(x, y) = 0 $.

Pentru a extrage funcția $ y = y(x) $, este necesar să punem condiția
de existență:
\[
  \dfrac{\partial F}{\partial y} \neq 0 \Rightarrow 3y^2 - 2x \neq 0.
\]

Acum, știind că avem funcția $ y = y(x) $, extremele acesteia se determină
folosind derivata $ y'(x) $, pe care nu o putem obține altfel decît prin
intermediul funcției $ F $. Avem, așadar:
\[
  \frac{\partial F}{\partial x} = 3x^2 - 3y^2 y' - 2y - 2x y' = 0.
\]
Obținem de aici:
\[
  y'(x) = \frac{2y - 3x^2}{3y^2 - 2x}.
\]
Pentru extreme, avem condițiile:
\[
  \begin{cases}
    y'(x) &= 0 \\
    F(x, y) &= 0 \\
    \dfrac{\partial F}{\partial y} &\neq 0
  \end{cases}
\]

Rezultă $ y = \dfrac{3x^2}{2} $ și înlocuim, obținînd perechea de soluții:
\[
  (x, y) = \Big( \dfrac{2\sqrt[3]{3}}{3}, \dfrac{2\sqrt[3]{4}}{3} \Big).
\]

Pentru a decide natura punctului de mai sus, este necesar să calculăm
$ y^{\prime\prime} $, care se obține derivînd din nou expresia pentru $ y'(x) $.
Avem:
\[
  y^{\prime\prime}(x) = \frac{(2y' - 6x)(3y^2 - 2x) - (6yy' - 2)(2y - 3x^2)}{(3y^2 - 2x)^2}.
\]
Evaluînd pentru $ x = \dfrac{2 \sqrt[3]{2}}{3} $, unde $ y' = 0 $, găsim:
\[
  y^{\prime\prime}\Big( \dfrac{2 \sqrt[3]{2}}{4} \Big) = -3 < 0,
\]
deci punctul este de maxim local.

%%%%%%%%%%%%%%%%%%%%%%%%%%%%%%%%%%%%%%%%%%%%%%%%%%%%%%%%%%%%%%%%%%%%%%

\section{Exerciții propuse}

\indent\indent 1. Fie funcția $ z = z(x, y) $, definită implicit prin:
\[
  g(y^2 - x^2, z - xy) = 0,
\]

unde $ g \in \kal{C}^1(\RR^2) $. Calculați expresia:
\[
  E = y \dfrac{\partial z}{\partial x} + x \dfrac{\partial z}{\partial y}.
\]

\vspace{0.5cm}

2. Să se determine extremele funcției $ y = y(x) $, definită implicit
de ecuațiile:
\begin{enumerate}[(a)]
\item $ x^3 + y^3 - 3x^2y - 3 = 0 $;
\item $ 2x^2y + y^2 - 4x - 3 = 0 $;
\item $ (x^2 + y^2)^2 = x^2 - y^2 $;
\item $ x^2 - 2xy + 5y^2 - 2x + 4y = - 1 $;
\item $ x^2 + y^2 - e^{2 \arctan \frac{x}{y}} = 0 $.
\end{enumerate}

\vspace{0.5cm}

3. Fie funcția $ z = z(x, y) $, definită implicit de ecuația:
\[
  F(x, y, z) = x^2 + 2y^2 + z^2 - 4z = 12.
\]
Să se calculeze derivatele parțiale de ordinul întîi și al doilea
pentru funcția $ z $.

%%% Local Variables:
%%% mode: latex
%%% TeX-master: "../bcd-19-20"
%%% End:
